\chapter{Giới thiệu}

\section{Bối cảnh và động lực}

\subsection{Bối cảnh}
Trong những năm gần đây, các mô hình học máy phức tạp như mạng nơ-ron sâu và phương pháp tập hợp cây đã được ứng dụng rộng rãi trong các quyết định quan trọng như chẩn đoán y tế, tuyển dụng và phê duyệt tín dụng. Mặc dù đạt độ chính xác dự đoán cao, các mô hình này thường thiếu tính giải thích, gây khó khăn cho việc hiểu và tin tưởng vào kết quả \parencite{doshivelez2017rigorousscienceinterpretablemachine}.

\subsection{Giải thích phản thực (Counterfactual Explanations)}
Một hướng tiếp cận phổ biến nhằm tăng tính minh bạch là các phương pháp giải thích cục bộ hậu nghiệm, đặc biệt là Giải thích phản thực (Counterfactual Explanations – CE) \parencite{wachter2018counterfactualexplanationsopeningblack}. CE đưa ra các hành động khả thi dưới dạng một vector nhiễu loạn giúp cá nhân thay đổi dự đoán sang kết quả mong muốn, ví dụ giảm chỉ số BMI để giảm nguy cơ tiểu đường. Quá trình này thường được mô hình hóa thành bài toán tối ưu, nhằm tìm hành động có chi phí thấp nhất nhưng vẫn đạt được kết quả mong muốn, còn được gọi là bài toán Giải pháp khắc phục khả thi (Actionable Recourse) \parencite{Ustun_2019}.

\subsection{Hạn chế của CE hiện tại}
Các phương pháp CE hiện nay chủ yếu tập trung vào từng trường hợp riêng lẻ. Trong thực tế, nhiều hành động có thể tác động đồng thời lên nhiều cá nhân, ví dụ khi doanh nghiệp áp dụng chính sách thăng chức hoặc tăng lương để giảm nguy cơ nghỉ việc của nhân viên. Trong bối cảnh này, việc gán hành động cần đáp ứng hai yêu cầu then chốt:

\begin{enumerate}
\item \textbf{Tính minh bạch:} Quy trình gán hành động phải được giải thích rõ ràng, tránh gây ra sự chênh lệch thiếu công bằng.
\item \textbf{Tính nhất quán:} Các lý do gán hành động phải thống nhất, không mâu thuẫn giữa các trường hợp tương tự.
\end{enumerate}

\begin{figure}[H]
\centering
\includegraphics[width=1\linewidth]{assets/Fig1.png}
\caption{Ví dụ về CET trên tập dữ liệu \textit{IBM HR Analytics Employee Attrition} \parencite{kaggle2017ibmhr}}
\label{fig:fig1}
\end{figure}

Các framework CE hiện hành và các phương pháp giải thích cục bộ khác chưa đáp ứng được hai yêu cầu này, do không xét đến toàn bộ không gian đầu vào \parencite{ribeiro2018anchors}.

\section{Mục tiêu nghiên cứu}
Để khắc phục hạn chế nêu trên, nghiên cứu này đề xuất \textbf{Cây Giải thích phản thực (Counterfactual Explanation Tree - CET)}. CET phân chia không gian đầu vào thành các miền con dễ hiểu, mỗi miền được gán một hành động thích hợp. Nhờ tính chất của cây quyết định, CET mang lại hai lợi thế:

\begin{enumerate}
\item \textbf{Tính minh bạch:} Quy tắc gán hành động được thể hiện dưới dạng cấu trúc cây dễ diễn giải, tương ứng với các luật if-then-else.
\item \textbf{Tính nhất quán:} Mỗi trường hợp chỉ thuộc một miền duy nhất, đảm bảo gán đúng một cặp hành động-lý do, tránh mâu thuẫn.
\end{enumerate}

\section{Đóng góp chính}
Nghiên cứu này đưa ra các đóng góp sau:

\begin{enumerate}
\item \textbf{Đề xuất CET:} Xây dựng cây quyết định gán hành động một cách minh bạch và nhất quán trên toàn bộ không gian đầu vào.
\item \textbf{Mô hình hóa và thuật toán:} Biểu diễn bài toán học CET dưới dạng tối ưu hóa và phát triển thuật toán tìm kiếm cục bộ ngẫu nhiên kết hợp chiến lược cắt tỉa dựa trên cận hàm mục tiêu.
\item \textbf{Thực nghiệm và nghiên cứu người dùng:} Thử nghiệm trên các bộ dữ liệu công khai cho thấy CET hiệu quả và dễ hiểu hơn so với các phương pháp hiện có. Nghiên cứu người dùng cũng chứng minh CET có khả năng gán cặp hành động–lý do rõ ràng và duy nhất.
\item \textbf{Tính linh hoạt:} CET có thể kết hợp với nhiều hàm chi phí khác nhau (ví dụ độ dịch chuyển bách phân vị lớn nhất \parencite{Ustun_2019}) và áp dụng cho cả các trường hợp chưa xuất hiện trong tập huấn luyện, tương tự các cơ chế giải thích phân bổ \parencite{chen2018learningexplaininformationtheoreticperspective}.
\end{enumerate}

\section{Cấu trúc báo cáo}
Báo cáo được tổ chức thành sáu chương:

\begin{itemize}[label={}]
\item \textbf{Chương 1: Giới thiệu} - Trình bày bối cảnh, động lực, mục tiêu và đóng góp chính của nghiên cứu.
\item \textbf{Chương 2: Các công trình liên quan} - Tổng quan các nghiên cứu về cây quyết định và Giải thích phản thực, chỉ ra khoảng trống nghiên cứu và vị trí của đề tài.
\item \textbf{Chương 3: Kiến thức nền tảng} - Trình bày các khái niệm cơ bản, cơ sở toán học và thuật toán được sử dụng.
\item \textbf{Chương 4: Phương pháp nghiên cứu} - Mô tả chi tiết phương pháp CET, thuật toán, mã giả và phân tích lý thuyết.
\item \textbf{Chương 5: Thực nghiệm và phân tích kết quả} - Trình bày thiết lập thực nghiệm, bộ dữ liệu, chỉ số đánh giá và phân tích kết quả.
\item \textbf{Chương 6: Kết luận và định hướng} - Tóm tắt đóng góp chính, hạn chế và hướng phát triển trong tương lai.
\end{itemize}
